\chapter{Preguntas frecuentes}
\label{sec:man-faq}

\section{¿Reemplaza a mi SIEM?}

No. El SIEM es mejor para búsquedas textuales por regla y para
alertas en tiempo real. GraphHunter es mejor para investigaciones
que requieren razonar sobre estructura (quién-autenticó-dónde-
conectó-adónde) sobre ventanas ya delimitadas. Los dos conviven.

\section{¿Cuándo me conviene activar scoring?}

Casi siempre. Los 5 componentes están calibrados para que un
grafo ``normal'' de oficina produzca distribuciones razonables.
Desactivarlo es útil si querés una pasada de exploración pura sin
priorización.

\section{¿Cuándo conviene activar el GNN threat?}

Sólo si tenés un modelo ONNX entrenado para tu entorno. Sin
modelo, \code{gnn_threat} devuelve 0 para todo y subir su peso
degrada el score total. Si no sabés, dejalo en 0.

\section{¿Puedo tener varias sesiones abiertas?}

Sí, en ventanas distintas del app. Cada sesión tiene su propio
grafo y su propio workspace. No se comparten tags ni notas entre
sesiones.

\section{¿El LLM ve mis datos?}

Sólo los que vos le entregás explícitamente en el panel agéntico
(muestras para proponer mappings). No pasa por él la ingesta
productiva. Si usás backend remoto (OpenAI/Anthropic/Google), esas
muestras salen por red; si usás backend local
(\code{GRAPHHUNTER_LOCAL_MODEL_PATH}), quedan en tu máquina.

\section{¿Qué pasa si ingesta un log corrupto?}

Los eventos que el parser entiende pero violan invariantes van a
la DLQ, no al grafo. Los eventos que ni siquiera parsean se loguean
como warning y se descartan, con conteo visible en el progreso de
ingesta.

\section{¿Cómo interpretar un path con gnn\_threat = 0?}

Significa que el GNN no está activo o no cargó modelo. No es un
score de amenaza \emph{baja}, es \emph{no hay señal}. El score
compuesto se apoya en los otros 4 componentes.

\section{¿Qué pasa si cambio el mapping sobre datos ya ingeridos?}

Los eventos ya están en el grafo con el mapping viejo. Cambiar el
mapping afecta sólo eventos futuros. Si querés reprocesar,
reingerí el archivo fuente.

\section{¿Puedo mover una sesión entre máquinas?}

Sí. \emph{Archivo \textrightarrow{} Exportar sesión} guarda un
archivo comprimido con el grafo, catálogo, hunts, notas y tags.
\emph{Importar sesión} lo restaura en la otra máquina. No se
exportan credenciales ni logs del sistema.

\section{¿Por qué el Host corporativo aparece con rarity alta?}

Probablemente tu ventana de ingesta es corta y ese host no tiene
mucha historia en el grafo. El \emph{entity rarity} se calcula
sobre el grafo actual, no sobre conocimiento externo. Es una
feature, no un bug: te obliga a pensar sobre qué entra al grafo.

\section{¿Puedo compartir un hunt con un colega?}

\begin{itemize}
    \item La hipótesis es texto DSL; copiar y pegar alcanza.
    \item Si querés los mismos resultados: exportar la sesión
    (archivo).
    \item Si querés reusar el scoring en una base compartida:
    publicar el mapping (si es custom) al \code{mapping_library}
    común.
\end{itemize}

\section{¿El grafo vive en disco?}

Mientras la sesión está abierta, vive en RAM para performance. Al
cerrar la sesión (o al exportar explícitamente), se persiste a
disco en el directorio de datos del app. Abrir una sesión
guardada la carga completa en RAM.

\section{¿Qué diferencia hay entre k, window y min\_score?}

\begin{itemize}
    \item \code{k}: simplicidad estructural. Cota de la explosión
    combinatoria del matcher.
    \item \code{window}: separación temporal máxima entre eventos
    del path.
    \item \code{min_score}: umbral \emph{post-hoc} sobre el score
    compuesto. No afecta el matching, sólo el filtrado de la
    salida.
\end{itemize}

\section{¿Por qué algunos endpoints requieren una flag de entorno?}

Acciones que publican fuera del app (IOCs a Sentinel, tickets a
Jira) están detrás de un switch explícito para que nadie las
dispare accidentalmente. La flag es la primera barrera; la
confirmación por llamada, la segunda.

\section{¿Mi hunt tarda mucho. Hay algo que pueda activar?}

Sí. El \emph{matcher} acepta varias \emph{flags} de entorno
\emph{opt-in} que activan rutas optimizadas incorporadas en la
revisión \emph{M-track} de abril 2026. Probar (en orden de
seguridad y \emph{lift}):

\begin{itemize}
    \item \code{GRAPHHUNTER_LFTJ_PLANNER=1}: si la hipótesis
    tiene \emph{shape} de \(K_4\) (cuatro entidades en clique con
    el mismo tipo y mismo tipo de relación entre todas),
    la dispatcha al enumerador WCOJ que es órdenes de magnitud
    más rápido.
    \item \code{GRAPHHUNTER_YANNAKAKIS=1}: \emph{matcher} de
    cadenas \emph{output-sensitive} con \emph{semi-join reducer}
    de dos pasadas. Para cadenas largas con buena selectividad
    elimina \emph{candidates} muertos antes de enumerar, así que
    rinde mejor que el DFS cuando la mayoría de los \emph{starts}
    no van a producir un \emph{match}.
    \item \code{GRAPHHUNTER_SEMIJOIN=1}: poda \emph{candidate
    sets} antes de entrar al DFS sobre cadenas heterogéneas.
    \item \code{GRAPHHUNTER_AMAC=1}: \emph{interleaved DFS} con
    \(K=8\) contextos en vuelo. Útil sobre todo en patrones de
    \(\geq 3\) saltos donde la latencia de DRAM domina.
    \item \code{GRAPHHUNTER_HUGEPAGES=1} (sólo Linux): reduce
    \emph{TLB stalls} en grafos grandes.
\end{itemize}

Todas son seguras: si la \emph{shape} no aplica o la flag no está,
se cae al \emph{path} legacy. Detalle completo en el libro
\emph{Especificación}, capítulo \emph{Variables de entorno}, y los
números medidos en \emph{Paper}, capítulo \emph{Evaluación}.

\section{¿Por qué mi sesión guardada cambió de formato?}

A partir de abril 2026, GraphHunter guarda las sesiones en un
formato \emph{binario} (\emph{magic bytes} \code{GHS1}) en lugar
del JSON anterior. Es transparente: las sesiones JSON existentes
siguen abriendo normalmente (el motor detecta el formato por los
primeros 4 bytes del archivo), pero al guardar de nuevo se
escriben en binario. Beneficios:

\begin{itemize}
    \item $\sim 5\times$ más chico en disco.
    \item $\sim 5\text{--}10\times$ más rápido para guardar y abrir.
    \item \emph{Streaming} en ambos sentidos: la RAM al cargar deja
    de ser proporcional al tamaño del log y queda acotada por el
    \emph{buffer} de I/O ($\sim$8 MB). Crítico para sesiones EVTX
    de 28 GB$+$ donde el JSON \emph{materializaba} todo el
    documento en memoria al cargar.
\end{itemize}

Si necesitás compartir una sesión binaria con una versión vieja
del cliente, abrila y volvé a guardarla con un cliente que todavía
escriba JSON --- ese caso es raro y la migración inversa no
está soportada.

\section{¿Por qué mi sesión grande abre rápido pero el primer hunt tarda más?}

A partir de abril 2026, los \emph{auxiliary indexes} (NLF y k-hop
reachability) se construyen \emph{lazy} --- al \emph{primer hunt}
en lugar de al cargar la sesión. Esto te permite ver el grafo y
navegarlo apenas terminó el \emph{streaming load}, en vez de
esperar 1--3 minutos a que se construyan los \emph{indexes} antes
de mostrarte la UI.

El primer \emph{hunt} sobre una sesión recién abierta paga el
costo del \emph{build} ($\sim$30~s para NLF, $\sim$1--2~min para
k-hop reachability en un EVTX de 28~GB). \emph{Hunts} subsiguientes
son inmediatos --- los \emph{indexes} están en memoria
(o en \emph{mmap-spilled tempfiles}) y se reutilizan.

Si querés evitar el primer-hunt-lento, dispará un \emph{hunt}
trivial apenas cargás la sesión: \code{User -[Auth]-> User} con
\emph{window} amplia es suficiente para forzar el \emph{build} en
\emph{background}.

\section{¿Cómo escribo una hipótesis con forma de estrella o árbol?}

A partir de abril 2026 la \emph{DSL} acepta el separador \code{;} para
introducir cláusulas adicionales que comparten \emph{bindings} con
cláusulas anteriores. Una estrella de 3 hojas centrada en \code{c}
se escribe:

\begin{verbatim}
User c -[Auth]-> Host h1;
c -[Auth]-> Host h2;
c -[Auth]-> Host h3
\end{verbatim}

El primer \code{;} marca la hipótesis como \emph{shape}-\code{Graph} y
el motor dispatcha automáticamente al \emph{tree-walk Yannakakis
enumerator}. Las cadenas lineales clásicas (sin \code{;}) siguen
parseándose exactamente como antes y dispatcheando al chain
matcher.

Después de un \code{;}, la cláusula puede arrancar de dos formas:

\begin{itemize}
    \item Una referencia \emph{bare} a un \emph{binding} declarado
    antes (como \code{c} arriba). El motor hereda el \emph{type} del
    \emph{binding} original.
    \item Una declaración fresca \code{Type binding} (e.g.,
    \code{; User u -[Auth]-> Host h}). Introduce un vértice
    independiente.
\end{itemize}

Si el \emph{binding} referenciado no existe en cláusulas anteriores,
el \emph{parser} retorna un error que lista los \emph{bindings}
válidos para que ajustes el nombre.

\section{¿Hay forma de correr varias hipótesis a la vez?}

Sí: el endpoint \emph{Tauri} \code{cmd_run_hunt_batch} acepta una
lista de hipótesis y las ejecuta secuencialmente sobre la misma
sesión, \emph{pre-warming} los índices auxiliares (NLF y k-hop
reachability) una vez antes de la primera consulta. Útil para:

\begin{itemize}
    \item Baterías de \emph{templates} ATT\&CK (decenas de
    \emph{patterns} sobre el mismo \emph{stream} EDR).
    \item Comparar varias hipótesis sobre la misma ventana
    temporal.
    \item Evitar el \emph{first-hunt-slow} en sesiones recién
    cargadas: la primera entrada del \emph{batch} paga el
    \emph{build} de los \emph{indexes} y el resto vuela.
\end{itemize}

Errores por hipótesis están aislados: una \emph{query} malformada
en la posición $k$ falla esa entrada con un \code{error}, pero el
resto del \emph{batch} sigue corriendo. La respuesta incluye un
flag \code{indexesWereWarm} que dice si la \emph{batch} encontró
los \emph{indexes} ya construidos (segunda \emph{batch} en
adelante) o si los tuvo que construir (primera \emph{batch} tras
\emph{open}).
